\section{Methods}
\label{sec:method}

\subsection{Data}

We use the TIGRIS corpus of 150 games scored on 32 axes~\cite{dellaLibera2026corpus}.
All axes are normalized to a 1--10 scale before analysis. Games with missing values
on more than 4 axes ($< 2\%$ of the corpus) are excluded from PCA analysis.

\subsection{Three Normalization Approaches}

We compare three normalization methods applied before PCA:

\textbf{Raw (N1):} No normalization beyond per-axis z-scoring across all games.
This is the standard approach and produces the dominant general quality PC.

\textbf{Weight-normalized (N2):} Each axis score is divided by BGG weight before
z-scoring, removing the portion of variance explained by overall game complexity.
This partially suppresses the general factor.

\textbf{Within-game z-score (N3):} For each game, we compute the mean and standard
deviation of that game's scores across all 32 axes and z-score each axis score
relative to that game's own distribution. This removes the game-level quality signal
entirely, leaving only the pattern of relative axis emphasis.

\subsection{PCA and Scree Analysis}

PCA is applied to each normalized dataset using singular value decomposition.
We report: scree plots showing variance explained per component, component loadings
for the top 5 PCs, and the top 3 axis correlates of each PC under N3.
Component labels are assigned post-hoc based on the highest-loading axes.
